% ===================================================================
% Front matter — Biographical Note
% ~0.5 page, Springer LNM / Cambridge convention.
% Monograph: Flos Aureus: A Theorem-Guided Architecture Search
%            through the Trinity Identity φ²+φ⁻²=3
% Author:    Кирилл Шамин (Cyril Shamin)
% Institute: University of Reading / gHashTag Open Doctorate
% ===================================================================
%
% TODO: This is a scaffold stub. The biographical text below uses a
%       placeholder identity. Replace with the actual author's bio
%       before submission:
%
%   TODO-1: Replace the name and affiliation with those of the
%           submitting candidate: Кирилл Шамин / Cyril Shamin,
%           University of Reading / gHashTag Open Doctorate.
%   TODO-2: Update Zenodo DOIs if new records are deposited between
%           now and submission.
%   TODO-3: Add the candidate's academic background (degrees held,
%           institutions, year of enrolment) — Cambridge §3.6
%           requires "name, degrees held, institution".
%   TODO-4: Keep total length ≤ 0.5 pages (Springer LNM §1.4).
%   TODO-5: List any peer-reviewed publications arising from the
%           thesis work (arXiv preprints are acceptable if no
%           peer-reviewed venue exists yet, but mark them clearly).
% ===================================================================

\chapter*{Biographical Note}
\addcontentsline{toc}{chapter}{Biographical Note}
\thispagestyle{plain}

\noindent
% TODO-REPLACE: Update the name and affiliations below.
\textbf{Cyril Shamin} (\textbf{Кирилл Шамин}) is a doctoral candidate
at the University of Reading and a member of the gHashTag Open
Doctorate Programme. His research sits at the intersection of formal
verification, algebraic combinatorics, and learned-system design,
unified by the \emph{Trinity Identity} \(\varphi^{2}+\varphi^{-2}=3\).

He is the originator of the Trinity Research Programme, which combines
mechanized proof in Coq with runtime invariants in Rust to bind
learned-system behaviour to formal algebraic identities over the
golden ratio. He is the author of the \textsf{trinity},
\textsf{trinity-clara}, and \textsf{trios} codebases under the
\texttt{gHashTag} organization, and the depositor of the Zenodo
records associated with the framework
(DOIs \href{https://doi.org/10.5281/zenodo.18947017}{18947017},
\href{https://doi.org/10.5281/zenodo.19227877}{19227877}, and
\href{https://doi.org/10.5281/zenodo.19227879}{19227879}).

His current research focuses on closing the remaining \emph{Admitted}
markers in the IGLA invariant set, on extending the empirical anchor
beyond CoBr\(_{2}\) and quasicrystals, and on connecting the Lucas
ring \(\mathcal{L} = \mathbb{Z}[\varphi]\) to higher-rank Coxeter
geometries.
